% ============================================================
%  Preamble -- Topics in Mathematics with Applications in Finance
% ============================================================
\usepackage[T1]{fontenc}
\usepackage[utf8]{inputenc}
\usepackage{amsmath,amsthm,mathtools}
\usepackage{newpxtext,newpxmath} % provides the AMS symbols: do not load amssymb
\usepackage{bm}
\usepackage{microtype}
% Screen layout (read on a laptop, not printed): two parallel columns (paracol), the text column
% with the width it had on A4 (15 cm) and a 6.5 cm right column for notes, answers, feedback,
% examples, proofs, course notes, history and asides. Every right-column item starts level with the
% paragraph it belongs to; if the right column is still busy, the text waits (white space).
% The old A4 print layout was: a4paper,inner=2.4cm,outer=3.6cm,marginparwidth=3.0cm,marginparsep=0.3cm (twoside).
\usepackage[paperwidth=25cm,paperheight=29.7cm,top=2.6cm,bottom=2.6cm,left=2cm,textwidth=22cm,
  marginparsep=0.5cm,marginparwidth=6.5cm,headheight=23pt]{geometry}
\usepackage{paracol}
\setcolumnwidth{15cm/0.5cm,6.5cm}
% counters shared by the two columns (equations in a proof, examples, answers ... in the right column)
% (sectioning counters cannot be global with this kernel: they are copied to the right column by
% \synccounter before each switch, see \side@emit)
\globalcounter{equation}\globalcounter{footnote}\globalcounter{figure}
\usepackage{booktabs,array,tabularx,longtable}
\usepackage[dvipsnames]{xcolor}
\usepackage[shortlabels,inline]{enumitem}
\usepackage{tikz}
\usetikzlibrary{arrows.meta,positioning,decorations.pathreplacing,calc}
\usepackage{pgfplots}
\pgfplotsset{compat=1.18}
\usepackage[most]{tcolorbox}
\usepackage{listings}
\usepackage{fancyhdr}
\usepackage[colorlinks=true,linkcolor=NavyBlue,urlcolor=NavyBlue,citecolor=NavyBlue]{hyperref}
\usepackage[nameinlink]{cleveref}

\setlist{itemsep=2pt,topsep=4pt}
\setlength{\parskip}{3pt}
\setlength{\emergencystretch}{2em}
\hfuzz=5pt

% ---------- headers ----------
\pagestyle{fancy}
\fancyhf{}
\fancyhead[L]{\small\leftmark}
\fancyhead[R]{\small\makebox[\marginparwidth][r]{\rightmark\quad\thepage}}
\renewcommand{\headrulewidth}{0.3pt}
\fancypagestyle{plain}{\fancyhf{}\fancyfoot[C]{\small\thepage}\renewcommand{\headrulewidth}{0pt}}

% ---------- colours (bright on purpose: yellow notes, blue definitions, red theorems, green examples) ----------
\definecolor{defcol}{HTML}{1565C0}  % blue: definitions
\definecolor{thmcol}{HTML}{D32F2F}  % red: theorems, propositions, lemmas, corollaries
\definecolor{intcol}{HTML}{2E7D32}  % green: examples (also the green of the figures)
\definecolor{notecol}{HTML}{FBC02D} % yellow: margin notes and their answers
\definecolor{intucol}{HTML}{8E24AA} % purple: intuition
\definecolor{fincol}{HTML}{EF6C00}  % orange: finance
\definecolor{cscol}{HTML}{00897B}   % teal: case studies
\definecolor{keycol}{HTML}{C2185B}  % magenta: key formulas
\definecolor{exercol}{HTML}{3949AB} % indigo: exercises
\definecolor{srccol}{HTML}{546E7A}  % grey-blue: sources, course notes, remarks
\definecolor{histcol}{HTML}{B8860B} % bronze: history boxes (classical columns)
\definecolor{doubtcol}{HTML}{C62828}
\definecolor{ideacol}{HTML}{1565C0}
\definecolor{okcol}{HTML}{388E3C}
\colorlet{anscol}{notecol!85!black}

% ---------- table rows: a little air between lines in every table ----------
\AddToHook{env/tabular/begin}{\renewcommand{\arraystretch}{1.35}}
\AddToHook{env/tabularx/begin}{\renewcommand{\arraystretch}{1.35}}
\AddToHook{env/longtable/begin}{\renewcommand{\arraystretch}{1.35}}

% ---------- theorem-like environments (numbered by chapter), coloured headings ----------
\newtheoremstyle{coldef}{}{}{\normalfont}{}{\bfseries\color{defcol}}{.}{ }{}
\newtheoremstyle{colex}{}{}{\normalfont}{}{\bfseries\color{intcol}}{.}{ }{}
\newtheoremstyle{colexer}{}{}{\normalfont}{}{\bfseries\color{exercol}}{.}{ }{}
\newtheoremstyle{colthm}{}{}{\itshape}{}{\bfseries\color{thmcol}}{.}{ }{}
\newtheoremstyle{colrem}{}{}{\normalfont}{}{\itshape\color{srccol}}{.}{ }{}
\theoremstyle{coldef}
\newtheorem{definition}{Definition}[chapter]
\theoremstyle{colex}
\newtheorem{example}[definition]{Example}
\theoremstyle{colexer}
\newtheorem{exercise}{Exercise}[chapter]
\theoremstyle{colthm}
\newtheorem{theorem}[definition]{Theorem}
\newtheorem{proposition}[definition]{Proposition}
\newtheorem{lemma}[definition]{Lemma}
\newtheorem{corollary}[definition]{Corollary}
\theoremstyle{colrem}
\newtheorem{remark}[definition]{Remark}
\crefname{exercise}{Exercise}{Exercises}
\globalcounter{definition}\globalcounter{exercise}

% Right-column items are queued and written into the right column at the next point where the
% text column is between paragraphs at top level (end of the paragraph, end of a box, end of a
% list): there the columns are synchronised, the queued items are typeset in the right column,
% and the text continues from where it was.
\makeatletter
\newif\ifnote@inbox
\newif\ifside@inright
\def\note@queue{}
\def\side@pcname{paracol}
\newcommand{\side@emit}{\ifx\note@queue\@empty\else
  \let\note@tmp\note@queue\gdef\note@queue{}%
  \synccounter{chapter}\synccounter{section}\synccounter{subsection}\synccounter{subsubsection}%
  \switchcolumn*\side@inrighttrue\note@tmp\par\side@inrightfalse\switchcolumn\fi}
\newcommand{\note@flush}{\ifvmode\ifnote@inbox\else\ifside@inright\else
  \ifx\@currenvir\side@pcname\side@emit\fi\fi\fi\fi}
\newcommand{\side@put}[1]{\gappto\note@queue{#1}\note@flush}
\AddToHook{para/after}{\note@flush}
\tcbset{note queue/.code={\appto\kvtcb@afterbox{\note@flush}}}
% every chapter/appendix file is one paracol environment: it starts right after the \chapter
% line of the file (\begin{paracol}{2}) and is closed here (appendix A, only long tables, has none:
% longtable does not work inside paracol)
\AddToHook{include/end}{\par\note@flush\ifx\@currenvir\side@pcname\end{paracol}\fi}
\newcommand{\flushside}{\par\note@flush}
\makeatother
\tcbset{thmbox/.style={enhanced,breakable,colback=#1!7,colframe=#1!7,borderline west={3pt}{0pt}{#1},
  sharp corners,left=6pt,right=4pt,top=2pt,bottom=2pt,before skip=8pt,after skip=8pt,
  before upper={\csname note@inboxtrue\endcsname},note queue}}
\tcolorboxenvironment{definition}{thmbox=defcol}
\tcolorboxenvironment{theorem}{thmbox=thmcol}
\tcolorboxenvironment{proposition}{thmbox=thmcol}
\tcolorboxenvironment{lemma}{thmbox=thmcol}
\tcolorboxenvironment{corollary}{thmbox=thmcol}
% (examples are boxed in the right column, see \begin{example} below)
\tcolorboxenvironment{exercise}{thmbox=exercol}
\tcolorboxenvironment{remark}{thmbox=srccol}

% ---------- content boxes ----------
\tcbset{notebox/.style={enhanced,breakable,colback=#1!7,colframe=#1,
  fonttitle=\bfseries,coltitle=white,colbacktitle=#1,boxrule=0.6pt,arc=2pt,
  left=6pt,right=6pt,top=4pt,bottom=4pt,before skip=10pt,after skip=10pt,
  before upper={\csname note@inboxtrue\endcsname},note queue}}
% Box title: the box type in capitals, then the topic
\newcommand{\boxtitle}[2]{\MakeUppercase{#1}\ifx&#2&\else: #2\fi}
% Intuition: the "physicist's reading" of a result
\newtcolorbox{intuition}[1][]{notebox=intucol,title={\boxtitle{Intuition}{#1}}}
% Finance: where the maths meets the market
\newtcolorbox{finance}[1][]{notebox=fincol,title={\boxtitle{In finance}{#1}}}
% Case study: summary of an R case study / notebook
\newtcolorbox{casestudy}[1][]{notebox=cscol,title={\boxtitle{Case study}{#1}}}
% Summary / key formulas
\newtcolorbox{keyformulas}{notebox=keycol,fontupper=\small,title={\boxtitle{Key formulas}{}}}
% Sources: which lectures, videos and files a chapter is built on
\newtcolorbox{sources}{notebox=srccol,fontupper=\footnotesize,title={\boxtitle{Sources for this chapter}{}},before upper=\raggedright}
\newcolumntype{L}{>{\raggedright\arraybackslash}X}
% Note on course material (errata, missing lectures, differences 2013/2024): right column, see below
% Secondary material ("beyond the core"): can be skipped on a first reading
\newtcolorbox{secondary}[1][]{enhanced,breakable,colback=white,colframe=white,boxrule=0pt,
  borderline west={1.5pt}{0pt}{srccol!60},sharp corners,left=8pt,right=0pt,top=1pt,bottom=1pt,
  before skip=6pt,after skip=6pt,fontupper=\small,
  title={\footnotesize\sffamily$\diamond$\ \boxtitle{Beyond the core}{#1}},
  coltitle=srccol,colbacktitle=white,fonttitle=\bfseries,titlerule=0pt,toptitle=0pt,bottomtitle=0pt,
  before upper={\csname note@inboxtrue\endcsname},note queue}
% ============================================================
%  The right column: the reader's notes/answers/feedback in \tiny; examples, proofs, course notes,
%  history and asides in \scriptsize
% ============================================================
%  \note{...}     the reader's comment/doubt, NEW (not yet read by Claude)
%  \note*{...}    a request that has been carried out (green DONE tag)
%  \note+{...}    a comment that has been read and needs no answer (blue SEEN tag)
%  \begin{answer}{question}[topic] ... \end{answer}
%                 a doubt and its answer in one box: the question on top, then a dark yellow
%                 bar "ANSWER n: topic" and the answer; \answerpart{topic} starts a further
%                 answer in the same box (one per question); \label{ans:...} makes it citable
%  \feedback{...} a comment on one of the reader's own solutions (indigo)
%  \begin{history}[topic], \begin{aside}[topic], \begin{coursenote}[topic]
%  example and proof environments are also printed here (sideexample = example)
%  Inside a coloured box of the main text, margin items appear right after the box.
\makeatletter
\tcbset{sidebox/.style={enhanced,colback=#1!25,colframe=#1,boxrule=0.5pt,arc=1pt,
    left=3pt,right=3pt,top=2pt,bottom=2pt,before skip=0pt,after skip=0pt,
    fontupper=\tiny\sffamily\raggedright,fonttitle=\tiny\sffamily\bfseries,
    before upper={\note@inboxtrue}}}
\newcommand{\side@tag}[2]{\colorbox{#1}{\textcolor{white}{\bfseries\tiny #2}}\par}
\NewDocumentCommand{\note}{s t+ +m}{\side@put{\begin{tcolorbox}[sidebox=notecol]%
  \IfBooleanT{#1}{\side@tag{okcol}{DONE}}\IfBooleanT{#2}{\side@tag{defcol}{SEEN}}#3\end{tcolorbox}}}
\newcommand{\feedback}[1]{\side@put{\begin{tcolorbox}[sidebox=exercol]\side@tag{exercol}{FEEDBACK}#1\end{tcolorbox}}}
\newcounter{answer}[chapter]
\renewcommand{\theanswer}{\thechapter.\arabic{answer}}
\crefname{answer}{answer}{answers}
\globalcounter{answer}
\newcommand{\answerpart}[1]{\refstepcounter{answer}%
  \tcbsubtitle[colback=anscol,colframe=anscol,coltext=black,fontupper=\tiny\sffamily\bfseries,
    left=3pt,top=0.5pt,bottom=0.5pt,before skip=2pt,after skip=2pt]{ANSWER \theanswer\ifx&#1&\else: #1\fi}}
\NewDocumentEnvironment{answer}{m O{} +b}{\side@put{\begin{tcolorbox}[sidebox=notecol,colback=notecol!12]%
  #1\answerpart{#2}#3\end{tcolorbox}}}{}
% Content boxes of the right column (examples, proofs, course notes, history, asides): \scriptsize,
% one step above the \tiny of the reader's notes
\tcbset{sidecontent/.style={sidebox=#1,colback=#1!6,breakable,
    fontupper=\scriptsize\raggedright,fonttitle=\scriptsize\sffamily\bfseries,coltitle=white,colbacktitle=#1}}
\NewDocumentEnvironment{aside}{O{} +b}{\side@put{\begin{tcolorbox}[sidecontent=histcol,
  title={\boxtitle{Aside}{#1}}]#2\end{tcolorbox}}}{}
\NewDocumentEnvironment{coursenote}{O{} +b}{\side@put{\begin{tcolorbox}[sidecontent=srccol,
  title={\boxtitle{Course note}{#1}}]#2\end{tcolorbox}}}{}
% Examples and proofs keep their theorem-like look (numbering, labels, QED) but go to the right column
\let\origexample\example \let\endorigexample\endexample
\RenewDocumentEnvironment{example}{o +b}{\side@put{\begin{tcolorbox}[thmbox=intcol,
  fontupper=\scriptsize,before skip=0pt,after skip=0pt]%
  \IfValueTF{#1}{\begin{origexample}[#1]}{\begin{origexample}}#2\end{origexample}\end{tcolorbox}}}{}
\NewDocumentEnvironment{sideexample}{o +b}{\IfValueTF{#1}{\begin{example}[#1]}{\begin{example}}#2\end{example}}{}
\let\origproof\proof \let\endorigproof\endproof
\RenewDocumentEnvironment{proof}{o +b}{\side@put{\begin{tcolorbox}[thmbox=thmcol,colback=thmcol!3,
  fontupper=\scriptsize,before skip=0pt,after skip=0pt]%
  \IfValueTF{#1}{\begin{origproof}[#1]}{\begin{origproof}}#2\end{origproof}\end{tcolorbox}}}{}
\makeatother
% History: a fact from the history of finance/mathematics, framed like a classical
% temple front: a dentilled entablature (the title) on two Ionic columns (not breakable)
\newcommand{\histcolumn}[2]{% #1 = foot of the column, #2 = top of the capital (coordinate names)
  \fill[histcol] ([xshift=-6pt]#1) rectangle ([xshift=6pt,yshift=2.5pt]#1);
  \fill[histcol] ([xshift=-4.5pt,yshift=2.5pt]#1) rectangle ([xshift=4.5pt,yshift=4pt]#1);
  \fill[histcol!15] ([xshift=-3.5pt,yshift=4pt]#1) rectangle ([xshift=3.5pt,yshift=-5pt]#2);
  \foreach \dx in {-3.5,-1.2,1.2,3.5}
    \draw[histcol,line width=0.4pt] ([xshift=\dx pt,yshift=4pt]#1) -- ([xshift=\dx pt,yshift=-5pt]#2);
  \fill[histcol] ([xshift=-4.5pt,yshift=-5pt]#2) rectangle ([xshift=4.5pt,yshift=-2.5pt]#2);
  \foreach \dx in {-4.5,4.5} {
    \draw[histcol,line width=0.8pt,fill=white] ([xshift=\dx pt,yshift=-4.5pt]#2) circle (1.9pt);
    \fill[histcol] ([xshift=\dx pt,yshift=-4.5pt]#2) circle (0.6pt);}
  \fill[histcol] ([xshift=-6.5pt,yshift=-2.5pt]#2) rectangle ([xshift=6.5pt]#2);}
\makeatletter
\NewDocumentEnvironment{history}{O{} +b}{\side@put{\begin{tcolorbox}[histbox,title={\boxtitle{History}{#1}}]#2\end{tcolorbox}}}{}
\makeatother
\tcbset{histbox/.style={enhanced,colback=histcol!6,colframe=histcol,boxrule=0.8pt,sharp corners,
  fonttitle=\scriptsize\sffamily\bfseries,fontupper=\scriptsize\raggedright,coltitle=white,colbacktitle=histcol,
  left=18pt,right=18pt,top=6pt,bottom=3pt,before skip=0pt,after skip=0pt,
  before upper={\csname note@inboxtrue\endcsname},
  overlay={
    \draw[histcol,line width=2.5pt,dash pattern=on 2.5pt off 2pt]
      ([yshift=-1.25pt]title.south west) -- ([yshift=-1.25pt]title.south east);
    \coordinate (hfL) at ([xshift=9pt]frame.south west);
    \coordinate (htL) at ([xshift=9pt,yshift=-2.5pt]title.south west);
    \coordinate (hfR) at ([xshift=-9pt]frame.south east);
    \coordinate (htR) at ([xshift=-9pt,yshift=-2.5pt]title.south east);
    \histcolumn{hfL}{htL}\histcolumn{hfR}{htR}}}}
% The reader's own solution to an exercise, placed right below the exercise
%   \begin{mysolution}[optional topic] ... \end{mysolution}
\newtcolorbox{mysolution}[1][]{notebox=exercol,colback=exercol!5,fontupper=\small,
  title={\boxtitle{My solution}{#1}}}
% A request about the notes themselves (collected in the "Requests" chapter)
%   \request{done|open|noted}{where it was written}{the request}{what was done}
\makeatletter
\newcommand{\request}[4]{\ifstrequal{#1}{done}{\def\req@col{okcol}}{\ifstrequal{#1}{open}{\def\req@col{thmcol}}{\def\req@col{defcol}}}%
  \begin{tcolorbox}[sidebox=notecol,colback=notecol!12,fontupper=\small\sffamily\raggedright,
  before skip=5pt,after skip=5pt]%
  \side@tag{\req@col}{\MakeUppercase{#1}}%
  {\itshape #2.}\ #3%
  \tcbsubtitle[colback=anscol,colframe=anscol,coltext=black,fontupper=\scriptsize\sffamily\bfseries,
    left=3pt,top=0.5pt,bottom=0.5pt,before skip=2pt,after skip=2pt]{WHAT WAS DONE}#4\end{tcolorbox}}
\makeatother



% ---------- video timestamps ----------
%  \ts{24}{12}{23:10}  -> lecture 12 of the 2024 course, at 23:10
%  \ts{13}{18}{5:02}   -> lecture 18 of the 2013 course
\makeatletter % consecutive tags may break between them, but never before punctuation
\DeclareRobustCommand{\ts}[3]{\mbox{\footnotesize\sffamily\color{srccol}[$\blacktriangleright$\,20#1\,L#2\,@\,#3]}\@ifnextchar\ts{\allowbreak}{}}
\makeatother

% ---------- listings (short R snippets) ----------
\lstset{language=R,basicstyle=\ttfamily\footnotesize,keywordstyle=\color{defcol}\bfseries,
  commentstyle=\color{intcol}\itshape,stringstyle=\color{thmcol},frame=single,rulecolor=\color{black!20},
  breaklines=true,columns=fullflexible,showstringspaces=false}

% ---------- maths shortcuts ----------
\newcommand{\R}{\mathbb{R}}
\newcommand{\N}{\mathbb{N}}
\newcommand{\E}{\mathbb{E}}
\newcommand{\Prob}{\mathbb{P}}
\newcommand{\Q}{\mathbb{Q}}
\DeclareMathOperator{\Var}{Var}
\DeclareMathOperator{\Cov}{Cov}
\DeclareMathOperator{\Corr}{Corr}
\DeclareMathOperator{\tr}{tr}
\DeclareMathOperator{\rank}{rank}
\DeclareMathOperator{\diag}{diag}
\newcommand{\T}{^{\mathsf T}}
\newcommand{\dd}{\,\mathrm{d}}
\newcommand{\ind}{\mathbf{1}}
\newcommand{\eqd}{\stackrel{d}{=}}

% ---------- file references ----------
\DeclareRobustCommand{\file}[1]{\texttt{\small\detokenize{#1}}}
